\documentclass[11pt,a4paper]{article}
\usepackage[utf8]{inputenc}
\usepackage[T1]{fontenc}
\usepackage{lmodern}
\usepackage[final]{microtype}
\usepackage{amsmath,amssymb,amsthm,mathtools,bm}
\usepackage{booktabs}
\usepackage{xcolor}
\usepackage[a4paper,top=28mm,bottom=30mm,left=30mm,right=30mm,headheight=14pt]{geometry}
\usepackage{fancyhdr}
\usepackage{titlesec}
\usepackage{enumitem}
\usepackage[font=small,labelfont=bf,margin=8mm]{caption}
\usepackage[pdftitle={A uniform quadratic wedge for consecutive Toeplitz minors of the Riemann xi-coefficients},pdfauthor={Alessio Doria - Independent Researcher},colorlinks=true,linkcolor=blue!50!black,citecolor=blue!50!black,urlcolor=blue!50!black]{hyperref}

\linespread{1.08}
\setlength{\parskip}{0.35em}
\setlist{itemsep=0.25em,topsep=0.4em}
\titleformat{\section}{\large\bfseries}{\thesection.}{0.6em}{}
\titleformat{\subsection}{\normalsize\bfseries}{\thesubsection.}{0.5em}{}
\titlespacing*{\section}{0pt}{1.6em}{0.7em}
\titlespacing*{\subsection}{0pt}{1.1em}{0.4em}
\numberwithin{equation}{section}

\pagestyle{fancy}
\fancyhf{}
\fancyhead[L]{\small\itshape A. Doria}
\fancyhead[R]{\small\itshape A uniform quadratic wedge for Toeplitz minors of $\xi$}
\fancyfoot[C]{\small\thepage}
\renewcommand{\headrulewidth}{0.4pt}
\fancypagestyle{plain}{\fancyhf{}\fancyfoot[C]{\small\thepage}\renewcommand{\headrulewidth}{0pt}}

\newtheorem{theorem}{Theorem}[section]
\newtheorem{proposition}[theorem]{Proposition}
\newtheorem{lemma}[theorem]{Lemma}
\newtheorem{corollary}[theorem]{Corollary}
\newtheorem*{mainthm}{Main Theorem}
\newtheorem*{inputA}{Input A (Michałowski, Gate A)}
\newtheorem*{inputB}{Input B (Liu--Mao, Lemma 3.3)}
\theoremstyle{definition}
\newtheorem{definition}[theorem]{Definition}
\theoremstyle{remark}
\newtheorem{remark}[theorem]{Remark}

\newcommand{\qq}{\mathfrak q}
\newcommand{\cA}{\mathcal A}
\newcommand{\cR}{\mathcal R}
\newcommand{\cI}{\mathcal I}
\newcommand{\cT}{\mathcal T}
\newcommand{\qbin}[3]{\genfrac{[}{]}{0pt}{}{#1}{#2}_{#3}}
\newcommand{\sgn}{\operatorname{sgn}}
\DeclareMathOperator{\diag}{diag}

\begin{document}

\thispagestyle{plain}
\begin{center}
{\LARGE\bfseries A uniform quadratic wedge for consecutive\\[0.35em] Toeplitz minors of the Riemann $\boldsymbol{\xi}$-coefficients\par}
\vspace{1.4em}
{\large Alessio Doria -- Independent Researcher\par}
\vspace{0.3em}
{\normalsize\href{mailto:numero2405@gmail.com}{\texttt{numero2405@gmail.com}}\par}
\vspace{0.6em}
{\normalsize November 2026\par}
\vspace{0.3em}
{\normalsize DOI: \href{https://doi.org/10.5281/zenodo.23946749}{10.5281/zenodo.23946749}\par}
\end{center}
\vspace{1.2em}

\begin{quote}\small
\noindent\textbf{Abstract.} Let $\tfrac18\xi(\tfrac12+\tfrac{ix}{2})=\sum_{k\ge0}(-1)^ka_kx^{2k}$, so that $a_k>0$, and let
$D_{r,k}=\det[a_{k+j-i}]_{i,j=0}^{r-1}$ be the consecutive Toeplitz minors of the coefficient sequence.
The Riemann Hypothesis implies that $(a_k)$ is a P\'olya frequency sequence, hence $D_{r,k}\ge0$ for all $r,k$.
Michałowski recently proved unconditionally that $D_{r,k}>0$ whenever $r\ge2$ and $k\ge10^{18}r^3$.
We prove that
\[
D_{r,k}>0\qquad\text{for all } r\ge2 \text{ and } k\ge 10^{11}r^2 .
\]
The improvement comes from replacing the three-point log-quadratic comparison sequence by a
four-point $q$-Pochhammer comparison sequence $c_s=(1-\alpha)^s/(\alpha;\qq)_s$ that reproduces the
normalized coefficients at four consecutive indices exactly. The reversed Hankel block of this reference is diagonally
congruent to a matrix with an exact $q$-Pascal $LDL^{T}$ factorization (a consequence of the
$q$-Chu--Vandermonde summation) with alternating pivots, so the dilation/whitening argument applies
to it unchanged, while the logarithmic mismatch now has four interpolation zeros. The weighted remainder
therefore starts at fourth order and scales like $r^2/k$ instead of $(r^3/k)^{1/2}$.
The only analytic inputs are Michałowski's all-degree Taylor bound for $\log a$ and the
Liu--Mao asymptotic for the log-curvature of $(a_k)$; all remaining constants are
certified in interval arithmetic.

\medskip
\noindent\textbf{Keywords.} Riemann $\xi$-function; Toeplitz minors; P\'olya frequency sequences; total positivity;
$q$-Pochhammer symbol; $q$-Chu--Vandermonde; Sylvester inertia.

\noindent\textbf{MSC 2020.} Primary 11M26; Secondary 15B05, 33D15, 30C15.
\end{quote}
\vspace{0.8em}


\section{Introduction}

\subsection{Setting}
Let $\xi(s)=\tfrac12 s(s-1)\pi^{-s/2}\Gamma(s/2)\zeta(s)$ and write, following \cite{Mich},
\begin{equation}\label{eq:Xi}
\Xi(x)=\frac18\,\xi\Big(\frac12+\frac{ix}{2}\Big)=\sum_{k\ge0}(-1)^ka_kx^{2k},\qquad
a_k=\frac{1}{(2k)!}\int_0^\infty u^{2k}\Phi(u)\,du>0,
\end{equation}
where $\Phi(u)=\sum_{n\ge1}(2\pi^2n^4e^{9u}-3\pi n^2e^{5u})\exp(-\pi n^2e^{4u})>0$.
For $r\ge1$ and $k\ge0$ set
\[
D_{r,k}=\det\big[a_{k+j-i}\big]_{i,j=0}^{r-1}\qquad(a_\ell=0\text{ for }\ell<0).
\]
The function $G(z)=\sum_k a_kz^k$ is entire of order $1/2$ and $\Xi(x)=G(-x^2)$. The Riemann Hypothesis (RH)
is equivalent to all zeros of $G$ being real and negative, and by the Aissen--Schoenberg--Whitney--Edrei theorem
\cite{ASW,Edrei} this happens if and only if $(a_k)$ is a P\'olya frequency sequence, i.e.\ every minor of the
Toeplitz matrix $[a_{j-i}]_{i,j\ge0}$ is nonnegative. The consecutive minors $D_{r,k}$ form a natural
two-parameter subfamily; under RH they are all nonnegative.

Unconditional information on $D_{r,k}$ is of several kinds. For orders two and three positivity is known for every $k$:
$D_{2,k}>0$ is the Turán inequality of Csordas, Norfolk and Varga \cite{CNV}, and $D_{3,k}>0$ is due to Nuttall
\cite{Nuttall} (see \cite{WY}). For each fixed $r$, Katkova \cite{Katkova} proved that $D_{r,k}>0$ for $k\ge N(r)$ with a
non-explicit $N(r)$. Wang and Yang report an order-uniform sufficient range of exponential size, $k>c\,e^{r(r-1)}$, for the
same determinants \cite[Introduction]{WY}. More recently Michałowski \cite{Mich} obtained the first explicit polynomial
uniform wedge,
\begin{equation}\label{eq:cubic}
D_{r,k}>0\qquad(r\ge2,\ k\ge10^{18}r^3).
\end{equation}

Several related results concern different hierarchies of the Laguerre--Pólya program. Wang and Yang \cite{WY} prove
Laguerre inequalities of order $r$ for $k>cr^3$ and complete-monotonicity estimates for the $\Xi$-coefficients;
Griffin--Ono--Rolen--Zagier \cite{GORZ} and Holland \cite{Holland} establish hyperbolicity regions for Jensen
polynomials of $\xi$; and Liu and Mao \cite{LM} prove that $(a_k)$ (equivalently $\lambda_n=\xi^{(2n)}(\tfrac12)/(2n)!$)
is strictly infinitely log-concave. Sra \cite{Sra} proves log-concavity and Pólya frequency properties of large
finite order for the densities whose Laplace transforms are $1/\xi(\tfrac12+s)$ and $1/\Xi$, which are different objects
from the coefficient sequence $(a_k)$. These conditions do not imply positivity of the order-$r$ consecutive Toeplitz minors
considered here. To the best of our knowledge, the result below is the first polynomial uniform wedge for these minors
with exponent strictly below three.

\subsection{Result}
\begin{mainthm}
For every integer $r\ge2$ and every integer $k\ge10^{11}r^2$ one has $D_{r,k}>0$.
\end{mainthm}

The constant $10^{11}$ is not optimized; the argument closes for every $C>1.35\cdot10^{10}$
(Remark~\ref{rem:constant}). The theorem says nothing about the RH-critical regime $k=O(r)$.

\subsection{Strategy}
We follow the architecture of \cite{Mich}: compare the reversed Hankel block
$H=[a_{k+d-i-j}]_{i,j=0}^{d}$, $d=r-1$, with an explicit reference block $M=SJS^T$ of known
inertia, and show that the \emph{whitened} perturbation $S^{-1}(H'-M)S^{-T}$ has spectral norm $<1$, where
$H'$ is a diagonal rescaling of $H$; Sylvester's law then fixes the sign of $\det H$.
The new ingredient is the reference. In \cite{Mich} the reference is $c_s=q_k^{s(s-1)/2}$, which matches
$\log a_{k+s}$ at the three points $s=-1,0,1$, so the logarithmic error $h$ has a cubic Taylor term of size
$k^{-2}$. Here we use the $q$-Pochhammer sequence
\[
c_s=\frac{(1-\alpha)^s}{(\alpha;\qq)_s},
\]
whose two parameters are chosen so that $c$ matches the normalized coefficients at $s=-1,0,1,2$. $LDL^T$ factorizations of $q$-Pascal type are classical; what is new here is their use for a four-point reference adapted to the $\xi$-coefficients. Two facts make
this work:
\begin{enumerate}
\item the reversed Hankel block of $c$ is diagonally congruent to the matrix
\[U_{ij}=\frac{(B;Q)_{i+j}}{(B;Q)_i\,(B;Q)_j},\]
which has an exact $LDL^T$ factorization with the \emph{same} $q$-Pascal factor $L=[\qbin ij{Q}]$ as the
$q$-Vandermonde matrix and with pivots of alternating sign (Proposition~\ref{prop:LDL});
\item the response of the whitened block to a diagonal dilation depends only on $L$ and on the pivot ratios, and the
new pivot ratios give a whitened sub-diagonal that is no larger than before (Lemma~\ref{lem:subdiag}).
\end{enumerate}
Since the error $h$ now vanishes at four nodes, it equals $\sum_{n\ge4}\beta_n(s^n-\cI_3[s^n])$
with $|\beta_n|\lesssim 40^nk^{1-n}$, and the relevant weight is $w\asymp\sqrt{rk}$; the leading term is
$k^{-3}w^4\asymp r^2/k$.

\subsection{Inputs and logical dependencies}
We use exactly two analytic results from the literature, stated in Section~\ref{sec:inputs}:
the $k$-unit form of Michałowski's all-degree bound (\cite[Prop.~3.7, first display]{Mich}) and
the Liu--Mao estimate \cite[Lemma~3.3]{LM}. We do \emph{not} use \cite[Lemma~2.1]{Mich} (the bounds
$1/(2k)<\tau_k<4/k$) nor its $\tau$-unit form of Gate A: all information on $\tau_k$ comes from Input~B.
The dilation estimates and the inertia lemma are re-proved here for the new reference.
All scalar inequalities are certified with Arb \cite{Arb}; scripts are provided as ancillary files.

\subsection*{Notation}
$(a;q)_n=\prod_{j=0}^{n-1}(1-aq^j)$ for $n\ge0$, $(a;q)_{-m}=1/(aq^{-m};q)_m$ for $m>0$, and
$\qbin nm{Q}=(Q;Q)_n/((Q;Q)_m(Q;Q)_{n-m})$. Throughout $r\ge2$, $d=r-1$, and $\|\cdot\|$ is the spectral norm.

\section{Analytic inputs}\label{sec:inputs}

Put $f=\log a$, where $a(z)=I(z)/\Gamma(2z+1)$ and $I(z)=\int_0^\infty u^{2z}\Phi(u)\,du$ is holomorphic on
$\operatorname{Re}z>-\tfrac12$; thus $a(n)=a_n$ for integers $n\ge0$.

\begin{inputA}[{\cite[Lemma~3.6, Prop.~3.7]{Mich}}]
For every integer $k\ge10^9$, $f$ is holomorphic on $|z-k|\le0.05k$ (with $f(n)=\log a_n$ at integers), and
\begin{equation}\label{eq:gateA}
\Big|\frac{f^{(n)}(k)}{n!}\Big|\le 3\cdot40^n\,k^{1-n}\qquad(n\ge3).
\end{equation}
\end{inputA}

Define the log-curvature
\[
q_k=\frac{a_{k-1}a_{k+1}}{a_k^2}=e^{-\tau_k},\qquad \tau_k>0 .
\]
Liu and Mao use $F(x)=\xi(\tfrac12+\sqrt x)=\sum\lambda_nx^n$. Since $\tfrac18\xi(\tfrac12+\tfrac{\sqrt z}{2})=\sum a_nz^n=\tfrac18F(z/4)$,
we have $a_n=\lambda_n/(8\cdot4^n)$, and therefore $\lambda_{n-1}\lambda_{n+1}/\lambda_n^2=q_n$.
Let $W$ be the principal branch of the Lambert function and
\[
\cT(z)=\frac{2\,W(2z/\pi)}{z\,(1+W(2z/\pi))}.
\]

\begin{inputB}[{\cite[Lemma~3.3, (3.53)]{LM}}]
For real $n\ge12000$ there is a function $g_0$, holomorphic on $D_0=\{z:|z-n|\le n/2\}$ and real on
$D_0\cap\mathbb R$, with $g_0(m)=\log(\lambda_m^2/\lambda_{m-1}\lambda_{m+1})=\tau_m$ for integers $m\in D_0$, and
\[
|g_0(z)-\cT(z)|\le\frac{20}{n^2}\qquad(z\in D_0).
\]
\end{inputB}

\begin{remark}
Input~A is used only through the Taylor coefficients of $f$ at an integer centre; any analytic interpolation of
$\log a_n$ with these bounds would do, and it need not coincide with the interpolation $g_0$ of Input~B, since
only the values at integers enter the comparison (Section~\ref{sec:remainder}).
\end{remark}

\section{Curvature windows}\label{sec:windows}

\begin{lemma}\label{lem:tau}
For every integer $k\ge12000$,
\[
1.71<\frac{12}{7}-\frac{20}{k}\le k\tau_k\le 2+\frac{20}{k}<2.002 .
\]
\end{lemma}

\begin{proof}
Write $w=W(2x/\pi)$. Since $6e^6<2\cdot12000/\pi$ and $W$ is increasing, $w>6$ for $x\ge12000$, so
$x\cT(x)=2w/(1+w)\in[12/7,2)$. Apply Input~B with $n=k$ at $z=k$.
\end{proof}

\begin{lemma}\label{lem:Delta}
For every integer $k\ge12000$, $\Delta_k:=\tau_k-\tau_{k+1}$ satisfies
\[
1.675<\frac{576}{343}\Big(\frac{k}{k+1}\Big)^2-\frac{40}{k-2}\le k^2\Delta_k\le 2+\frac{40}{k-2}<2.004 .
\]
In particular $q_k<q_{k+1}$.
\end{lemma}

\begin{proof}
With $w=W(2x/\pi)$ and $w'(x)=w/(x(1+w))$ one computes
\[
-\cT'(x)=\frac{2}{x^2}\,g(w),\qquad g(w)=\frac{w^2(w+2)}{(1+w)^3}.
\]
Since $(\log g)'(w)=\frac2w+\frac1{w+2}-\frac3{w+1}=\frac{w+4}{w(w+1)(w+2)}>0$, $g$ increases from
$g(6)=288/343$ to $1$; hence $\frac{576}{343}x^{-2}\le-\cT'(x)<2x^{-2}$ for $x\ge12000$, and
$\cT(k)-\cT(k+1)\in[\frac{576}{343}(k+1)^{-2},\,2k^{-2}]$.
Let $E=g_0-\cT$ with $n=k$ in Input~B. For $\eta\in[k,k+1]$ the disc $|z-\eta|\le k/2-1$ lies in $D_0$, so
Cauchy's estimate gives $|E'(\eta)|\le(20/k^2)/(k/2-1)=40/(k^2(k-2))$, whence
$|E(k)-E(k+1)|\le40/(k^2(k-2))$. Both bounds are monotone in $k$ in the favourable direction; the numerical
values at $k=12000$ are certified in Appendix~\ref{app:cert}.
\end{proof}

\section{The four-point \texorpdfstring{$q$}{q}-Pochhammer reference}\label{sec:reference}

Fix an integer $k\ge12000$ and normalize
\[
\rho=\frac{a_{k+1}}{a_k},\qquad b_s=\frac{a_{k+s}}{a_k\,\rho^s},
\]
so that $b_{-1}=q_k=:u$, $b_0=b_1=1$, $b_2=q_{k+1}=:v$, and $0<u<v<1$ by Lemmas~\ref{lem:tau}--\ref{lem:Delta}.

\begin{definition}
Let
\[
\qq=\frac{1-v}{v(1-u)},\qquad \alpha=\frac{(1-u)(1-v)}{v-u},\qquad \theta=\log\qq,\qquad
c_s=\frac{(1-\alpha)^s}{(\alpha;\qq)_s}\quad(s\in\mathbb Z).
\]
\end{definition}

\begin{lemma}\label{lem:match}
$c_{-1}=u$, $c_0=c_1=1$, $c_2=v$. Moreover, with $\sigma=\tau_{k+1}$ and $\Delta=\tau_k-\sigma$,
\[
\alpha=\frac{2\sinh(\tau_k/2)\sinh(\sigma/2)}{\sinh(\Delta/2)},\qquad
\theta=\frac{\tau_k+\sigma}{2}+\log\frac{\sinh(\sigma/2)}{\sinh(\tau_k/2)} .
\]
\end{lemma}

\begin{proof}
$c_0=1$, $c_1=1$, $c_2=(1-\alpha)/(1-\alpha\qq)$ and $c_{-1}=(1-\alpha\qq^{-1})/(1-\alpha)$. Substituting the
definitions of $\alpha,\qq$ gives $c_2=v$ and $c_{-1}=u$ after clearing denominators (a polynomial identity in
$u,v$). The hyperbolic forms follow from $1-e^{-x}=2e^{-x/2}\sinh(x/2)$.
\end{proof}

\begin{lemma}\label{lem:params}
For $k\ge12000$: $1.45<\alpha<2.40$ and $0.53<k\theta<2.002$. In particular $\qq>1$.
\end{lemma}

\begin{proof}
Using $x\le\sinh x\le xe^{x^2/6}$,
$\alpha$ lies between $\frac{\tau_k\sigma}{\Delta}e^{-\Delta^2/24}$ and $\frac{\tau_k\sigma}{\Delta}e^{(\tau_k^2+\sigma^2)/24}$.
For $\theta$: $(\tau_k+\sigma)/2=\sigma+\Delta/2\ge\sigma$, $\log(\sigma/\tau_k)=-\log(1+\Delta/\sigma)\ge-\Delta/\sigma$,
and $\log\frac{\sinh(\sigma/2)}{\sinh(\tau_k/2)}\ge\log\frac{\sigma}{\tau_k}-\frac{\tau_k^2}{24}$; hence
$\theta\ge\sigma-\Delta/\sigma-\tau_k^2/24$. Also $\theta<\tau_k$ because $\sinh$ is increasing. Insert
Lemmas~\ref{lem:tau}--\ref{lem:Delta} and $k\sigma\ge k\tau_k-k\Delta_k$; see Appendix~\ref{app:cert}.
\end{proof}

\begin{lemma}[Positivity]\label{lem:positivity}
If $(r-1)\theta<\log\alpha$, then $c_s>0$ for $|s|\le r-1$, and every factor $1-\alpha\qq^{j}$ with
$-(r-1)\le j\le r-2$ is negative.
\end{lemma}

\begin{proof}
$\alpha\qq^{j}\ge\alpha e^{-(r-1)\theta}>1$ for $j\ge-(r-1)$. For $s>0$ numerator and denominator of $c_s$ both have sign
$(-1)^s$; for $s=-m<0$, $c_s=(1-\alpha)^{-m}\prod_{j=1}^m(1-\alpha\qq^{-j})$ has $2m$ negative factors.
\end{proof}

\begin{remark}
The condition $1<\qq<\alpha$ alone does \emph{not} suffice for positivity at negative indices; the hypothesis
$(r-1)\theta<\log\alpha$ is what is needed, and it holds in the wedge of the Main Theorem with enormous room
(Appendix~\ref{app:cert}).
\end{remark}

\subsection{Analytic continuation and Taylor coefficients}
Put $z_0=1/\alpha\in(0.41,0.69)$ and $Q=\qq^{-1}=e^{-\theta}$. Writing $1-\alpha\qq^j=-\alpha\qq^j(1-z_0Q^j)$ gives
$c_s=(1-z_0)^sQ^{s(s-1)/2}/(z_0;Q)_s$ for integers $s$ (for $s<0$ by the cocycle relation
$(a;q)_{m+n}=(a;q)_m(aq^m;q)_n$, valid for all integers $m,n$). With $(z_0;Q)_s=(z_0;Q)_\infty/(z_0Q^s;Q)_\infty$
and $\log(x;Q)_\infty=-\sum_{m\ge1}x^m/(m(1-Q^m))$ we obtain, for $|s|<\theta^{-1}\log(1/z_0)$,
\begin{equation}\label{eq:logc}
\ell(s):=s\log(1-z_0)-\frac\theta2 s(s-1)+\sum_{m\ge1}\frac{z_0^m\big(1-e^{-m\theta s}\big)}{m\,(1-e^{-m\theta})},
\end{equation}
a holomorphic function, real on the real axis, with $e^{\ell(s)}=c_s$ at integers $s$ in its domain.

\begin{lemma}\label{lem:refcoef}
For $n\ge3$, $|[s^n]\,\ell(s)|\le\kappa^n/n+\kappa^{n-1}/(n(n-1))$ with $\kappa=\theta/(1-z_0)$. For $k\ge12000$
this gives $|[s^n]\ell(s)|\le(9/k)^{n-1}$.
\end{lemma}

\begin{proof}
For $n\ge3$ only the series contributes, and $[s^n](1-e^{-m\theta s})=-(-m\theta)^n/n!$, so
\[
|[s^n]\ell|\le\frac{\theta^n}{n!}\sum_{m\ge1}\frac{m^{n-1}z_0^m}{1-e^{-m\theta}}
\le\frac{\theta^n}{n!}\Big(\sum_m m^{n-1}z_0^m+\frac1\theta\sum_m m^{n-2}z_0^m\Big),
\]
using $1/(1-e^{-x})\le1+1/x$. Since $\sum_{m\ge1}m^pz^m=zA_p(z)/(1-z)^{p+1}$ with the Eulerian polynomial
$A_p$ (nonnegative coefficients, $A_p(1)=p!$), $\sum_m m^pz^m\le p!/(1-z)^{p+1}$ for $0<z<1$, $p\ge0$. This gives the
first bound. By Lemma~\ref{lem:params}, $k\kappa=k\theta\,\alpha/(\alpha-1)<2.002\cdot1.45/0.45<6.46<8.4$, and
$8.4^n/(nk)+8.4^{n-1}/(n(n-1))<9^{n-1}$ for $n\ge3$, $k\ge12000$.
\end{proof}

\section{The reversed Hankel block and its \texorpdfstring{$q$}{q}-Pascal factorization}

Set $d=r-1$ and
\[
M_{ij}=c_{d-i-j},\qquad 0\le i,j\le d,\qquad Q=\qq^{-1},\qquad B=\alpha\qq^{d-1}.
\]

\begin{lemma}[Congruence]\label{lem:congruence}
$M=c_d\,A\,U\,A$ with $A=\diag(A_i)$, $A_i=(B;Q)_i/(1-\alpha)^i$, and
\[
U_{ij}=\frac{(B;Q)_{i+j}}{(B;Q)_i\,(B;Q)_j}.
\]
Under the hypothesis of Lemma~\ref{lem:positivity}, $c_d>0$ and $A_i>0$.
\end{lemma}

\begin{proof}
By the integer cocycle, $(\alpha;\qq)_d=(\alpha;\qq)_{d-n}(\alpha\qq^{d-n};\qq)_n$ for every $0\le n\le 2d$, including
$d-n<0$, and $(\alpha\qq^{d-n};\qq)_n=\prod_{l=0}^{n-1}(1-\alpha\qq^{d-1-l})=(B;Q)_n$. Hence
$c_{d-n}=c_d\,(B;Q)_n/(1-\alpha)^n$, which is the claim with $n=i+j$. Signs follow from Lemma~\ref{lem:positivity}:
$1-BQ^l=1-\alpha\qq^{d-1-l}<0$ for $0\le l\le d-1$, so $(B;Q)_i$ and $(1-\alpha)^i$ both have sign $(-1)^i$.
\end{proof}

\begin{proposition}[$q$-Pascal $LDL^T$]\label{prop:LDL}
For indeterminates $B,Q$ and all $i,j\ge0$,
\[
U_{ij}=\sum_{m=0}^{\min(i,j)}\qbin im{Q}\qbin jm{Q}D_m,\qquad
D_m=B^mQ^{m(m-1)}\frac{(Q;Q)_m}{(B;Q)_m}.
\]
Thus $U=LDL^T$ with $L_{ij}=\qbin ij{Q}$ and $D=\diag(D_m)$. In our situation $0<Q<1$ and $BQ^{m-1}=\alpha\qq^{d-m}>1$
for $1\le m\le d$, so $\sgn D_m=(-1)^m$.
\end{proposition}

\begin{proof}
From $(Q;Q)_i/(Q;Q)_{i-m}=(-1)^mQ^{im-\binom m2}(Q^{-i};Q)_m$ we get
\[
\qbin im{Q}\qbin jm{Q}D_m=\frac{(Q^{-i};Q)_m(Q^{-j};Q)_m}{(Q;Q)_m(B;Q)_m}\,(BQ^{i+j})^m .
\]
Hence the sum equals ${}_2\phi_1(Q^{-i},Q^{-j};B;Q,BQ^{i+j})$, which by the $q$-Chu--Vandermonde summation
${}_2\phi_1(q^{-n},b;c;q,cq^n/b)=(c/b;q)_n/(c;q)_n$ \cite[\S1.5 and Appendix~II]{GR}, with $n=i$, $b=Q^{-j}$, $c=B$, equals
$(BQ^j;Q)_i/(B;Q)_i=U_{ij}$. The sign of $D_m$: $(Q;Q)_m>0$ and every factor $1-BQ^l$, $0\le l\le m-1$, of $(B;Q)_m$ is negative.
\end{proof}

Consequently
\begin{equation}\label{eq:SJS}
M=SJS^T,\qquad S=\sqrt{c_d}\;A\,L\,|D|^{1/2},\qquad J=\diag\big((-1)^m\big)_{m=0}^d .
\end{equation}

\section{Dilation response of the whitened reference}\label{sec:gateB}

For a function $\phi$ defined at the integers $|s|\le d$ write $M\circ\phi:=[c_{d-i-j}\phi(d-i-j)]_{i,j}$ and define the
\emph{response}
\[
\cR(\phi)=\big\|S^{-1}(M\circ\phi)S^{-T}\big\| .
\]
Since $M\circ\phi=c_dA(U\circ\phi)A$ and diagonal matrices commute with the Hadamard product structure,
$\cR(\phi)=\||D|^{-1/2}L^{-1}(U\circ\phi)L^{-T}|D|^{-1/2}\|$; in particular $A$ and $c_d$ play no role.

\begin{lemma}\label{lem:R1}
Let $\Lambda_\beta=\diag(Q^{\beta i})_{i=0}^d$ and $R_\beta=L^{-1}\Lambda_\beta L$. Then $R_1$ is lower bidiagonal with
\[
(R_1)_{mm}=Q^m,\qquad (R_1)_{m,m-1}=-Q^{m-1}(1-Q^m).
\]
\end{lemma}

\begin{proof}
We need $\Lambda_1L=LR_1$, i.e.\ $Q^i\qbin ij{Q}=Q^j\qbin ij{Q}-Q^j(1-Q^{j+1})\qbin i{j+1}{Q}$. This follows from the
$q$-Pascal relation $(1-Q^{j+1})\qbin i{j+1}{Q}=(1-Q^{i-j})\qbin ij{Q}$.
\end{proof}

\begin{lemma}[Whitened sub-diagonal]\label{lem:subdiag}
Let $\widehat R_\beta=|D|^{-1/2}R_\beta|D|^{1/2}$. Then $(\widehat R_1)_{mm}=Q^m$ and
\[
\big|(\widehat R_1)_{m,m-1}\big|=\sqrt{(1-Q^m)\big(Q^{m-1}-B^{-1}\big)}\le\sqrt{1-Q^m}\le\sqrt{m\theta}.
\]
\end{lemma}

\begin{proof}
By Proposition~\ref{prop:LDL},
\[
\frac{D_m}{D_{m-1}}=-\frac{BQ^{2m-2}(1-Q^m)}{BQ^{m-1}-1},\qquad
(\widehat R_1)_{m,m-1}=(R_1)_{m,m-1}\,\frac{|D_{m-1}|^{1/2}}{|D_m|^{1/2}}.
\]
Squaring and simplifying gives the identity; then use
$0<Q^{m-1}-B^{-1}<1$ and $1-e^{-m\theta}\le m\theta$.
\end{proof}

\begin{remark}
For the $q$-Vandermonde reference of \cite{Mich} the corresponding quantity is $\sqrt{1-Q^m}$; the new reference
multiplies it by $\sqrt{Q^{m-1}-B^{-1}}\approx\sqrt{1-1/\alpha}<1$. Nothing else in the dilation argument depends on the
pivots.
\end{remark}

\begin{lemma}[Generator]\label{lem:generator}
Let $x=r\theta\le1/64$. Then $N:=\widehat R_1-I$ satisfies $\|N\|\le x+\sqrt x\le 9/64$, the principal logarithm
$G=\log\widehat R_1=\sum_{j\ge1}(-1)^{j+1}N^j/j$ satisfies $\|G\|\le\frac{72}{55}\sqrt x$, $\widehat R_\beta=e^{\beta G}$
for all real $\beta$, and
\[
\theta d+2\|G\|\le4\sqrt{r\theta}=:t .
\]
\end{lemma}

\begin{proof}
The diagonal part of $N$ has entries $|Q^m-1|\le m\theta\le x$ and its sub-diagonal part has entries $\le\sqrt x$
(Lemma~\ref{lem:subdiag}); a diagonal and a weighted shift have norms equal to their largest entries, so
$\|N\|\le x+\sqrt x$. Then $\|G\|\le-\log(1-\|N\|)\le\|N\|/(1-\|N\|)\le\frac{64}{55}\cdot\frac98\sqrt x$.
Next, $R_\beta=\exp(\beta L^{-1}\Lambda' L)$ with $\Lambda'=\diag(i\log Q)$, so $\widehat R_\beta=e^{\beta\widehat G}$ with
$\widehat G=|D|^{-1/2}L^{-1}\Lambda'L|D|^{1/2}$. Both $\widehat G$ and $G$ equal $P(\widehat R_1)$, where $P$ is the
Lagrange polynomial with $P(Q^i)=i\log Q$ (the eigenvalues $Q^i$ of $\widehat R_1$ are distinct and $\widehat R_1$ is
diagonalizable); hence $G=\widehat G$. Finally $\theta d+2\|G\|\le x+\frac{144}{55}\sqrt x\le(\frac18+\frac{144}{55})\sqrt x<4\sqrt x$.
\end{proof}

\begin{proposition}[All-degree response]\label{prop:response}
If $r\theta\le1/64$, then $\cR\big((\theta s)^n\big)\le t^n$ for every integer $n\ge0$, with $t=4\sqrt{r\theta}$.
\end{proposition}

\begin{proof}
For $s=d-i-j$ we have $e^{\beta\theta s}=e^{\beta\theta d}Q^{\beta i}Q^{\beta j}$, so
$M\circ e^{\beta\theta s}=e^{\beta\theta d}\Lambda_\beta M\Lambda_\beta$. Since $A$ commutes with $\Lambda_\beta$,
$S^{-1}\Lambda_\beta S=|D|^{-1/2}L^{-1}\Lambda_\beta L|D|^{1/2}=\widehat R_\beta=e^{\beta G}$, and therefore
\[
F(\beta):=S^{-1}\big(M\circ e^{\beta\theta s}\big)S^{-T}=e^{\beta\theta d}\,e^{\beta G}\,J\,e^{\beta G^T}.
\]
Entrywise, $F^{(n)}(0)=S^{-1}(M\circ(\theta s)^n)S^{-T}$. Expanding the three exponentials,
$F^{(n)}(0)=\sum_{a+b+c=n}\frac{n!}{a!\,b!\,c!}(\theta d)^aG^bJ(G^T)^c$, whose norm is at most $(\theta d+2\|G\|)^n\le t^n$.
\end{proof}

\begin{definition}[Weighted algebra]
For $w\ge1$ let $\cA_w$ be the space of power series $\phi(s)=\sum_{n\ge0}\phi_ns^n$ with
$\|\phi\|_{\cA}=\sum_n|\phi_n|w^n<\infty$. It is a commutative Banach algebra, and
$\|e^\phi-1\|_{\cA}\le e^{\|\phi\|_\cA}-1$.
\end{definition}

\begin{corollary}\label{cor:calculus}
Let $r\theta\le1/64$ and $w=t/\theta=4\sqrt{r/\theta}$. Then $w\ge32r>d$, and for every $\phi\in\cA_w$,
$\cR(\phi)\le\|\phi\|_{\cA}$.
\end{corollary}

\begin{proof}
$w/r=4/\sqrt{r\theta}\ge32$. The series $\phi(s)=\sum\phi_ns^n$ converges absolutely at every integer $|s|\le d$, so
$M\circ\phi=\sum_n\phi_n\,M\circ s^n$ entrywise, and $\cR(s^n)=\theta^{-n}\cR((\theta s)^n)\le w^n$.
\end{proof}

\section{The four-zero remainder}\label{sec:remainder}

Let $k\ge10^9$ and $r\theta\le1/64$. By Input~A, $\log b_s=f(k+s)-f(k)-s\,(f(k+1)-f(k))$ at integers $|s|\le d$, and the
right-hand side is holomorphic for $|s|<0.05k$. Define
\[
h(s)=f(k+s)-f(k)-s\,(f(k+1)-f(k))-\ell(s),\qquad \beta_n=[s^n]h(s).
\]
Then $b_s=c_s\,e^{h(s)}$ at the integers $|s|\le d$, and, by Lemma~\ref{lem:match},
\[
h(-1)=h(0)=h(1)=h(2)=0 .
\]

\begin{lemma}\label{lem:beta}
For $n\ge3$, $|\beta_n|\le3\cdot40^nk^{1-n}+(9/k)^{n-1}\le4\cdot40^nk^{1-n}$.
\end{lemma}

\begin{proof}
For $n\ge2$, $\beta_n=f^{(n)}(k)/n!-[s^n]\ell$; use \eqref{eq:gateA}, Lemma~\ref{lem:refcoef} and $9^{n-1}\le40^n$.
\end{proof}

Let $\cI_3$ denote Lagrange interpolation at the nodes $-1,0,1,2$, with basis
\[
\ell_{-1}(s)=-\tfrac{s^3}6+\tfrac{s^2}2-\tfrac s3,\quad \ell_0(s)=\tfrac{(s+1)(s-1)(s-2)}{2},\quad
\ell_1(s)=-\tfrac{s^3}2+\tfrac{s^2}2+s,\quad \ell_2(s)=\tfrac{s^3}6-\tfrac s6,
\]
and put $P_n(s)=s^n-\cI_3[s^n]$; note $P_n=0$ for $n\le3$.

\begin{lemma}\label{lem:interp}
If $w\ge1$ and $n\ge4$, then $\|\cI_3[s^n]\|_\cA\le2^nw^3$ and $\|P_n\|_\cA\le w^n+2^nw^3$.
\end{lemma}

\begin{proof}
For $n\ge1$, $\cI_3[s^n]=(-1)^n\ell_{-1}+\ell_1+2^n\ell_2$. For $w\ge1$, $\|\ell_{-1}\|_\cA\le w^3$,
$\|\ell_1\|_\cA\le2w^3$, $\|\ell_2\|_\cA\le w^3/3$, and $3+2^n/3\le2^n$ for $n\ge3$.
\end{proof}

\begin{proposition}\label{prop:majorant}
Assume $k\ge10^9$, $r\theta\le1/64$, $w=4\sqrt{r/\theta}$ and $y:=40w/k<1$. Then
\[
h=\sum_{n\ge4}\beta_nP_n\quad\text{in }\cA_w,\qquad
\|h\|_\cA\le 4k\,\frac{y^4}{1-y}+4k\,w^3\,\frac{(80/k)^4}{1-80/k}.
\]
\end{proposition}

\begin{proof}
By Lemmas~\ref{lem:beta} and~\ref{lem:interp}, the series $\sum_{n\ge4}\beta_nP_n$ converges absolutely
in $\cA_w$, and $\sum_n\beta_n\cI_3[s^n]$ converges to a cubic polynomial. Hence
\[
p:=h-\sum_{n\ge4}\beta_nP_n=\sum_{n\le3}\beta_ns^n+\sum_{n\ge4}\beta_n\cI_3[s^n]
\]
is a cubic polynomial; it vanishes at the
four nodes because $h$ and every $P_n$ do, so $p\equiv0$. The bound follows by summing
$4\cdot40^nk^{1-n}(w^n+2^nw^3)$ over $n\ge4$.
\end{proof}

The first term is the quartic main term: since $w^2=16r/\theta$ and $k\theta\asymp1$, it is $\asymp k^{-3}w^4\asymp r^2/k$.
For the three-point reference the analogous main term is $\asymp k^{-2}w^3\asymp(r^3/k)^{1/2}$, which is the source of the
cubic wedge \eqref{eq:cubic}.

\section{Proof of the Main Theorem}\label{sec:proof}

\begin{lemma}[Inertia]\label{lem:inertia}
Let $S$ be invertible, $J$ a signature matrix, and $T=S(J+F)S^T$ with $F$ symmetric and $\|F\|<1$. Then $T$ and $J$ have
the same inertia.
\end{lemma}

\begin{proof}
For $\lambda\in[0,1]$, $\sigma_{\min}(J+\lambda F)\ge1-\lambda\|F\|>0$, so the eigenvalues of the symmetric family
$J+\lambda F$ never cross $0$; hence $J+F$ has the inertia of $J$, and Sylvester's law of inertia transfers it to $T$.
\end{proof}

\begin{proof}[Proof of the Main Theorem]
Let $r\ge2$ and $k\ge10^{11}r^2$; then $k\ge4\cdot10^{11}$, so Inputs~A and~B and all lemmas above apply.
By Lemma~\ref{lem:params}, $r\theta<2.002\,r/k\le2.002/(10^{11}r)$, which is $\le1/64$ and $<\log1.45<\log\alpha$;
so Lemma~\ref{lem:positivity}, \eqref{eq:SJS}, Corollary~\ref{cor:calculus} and Proposition~\ref{prop:majorant} hold.
Since $k\theta>0.53$, $w^2=16r/\theta<16rk/0.53$ and
\[
y^2=\frac{1600\,w^2}{k^2}<\frac{25600}{0.53}\cdot\frac rk\le\frac{25600}{0.53\cdot2\cdot10^{11}}<2.5\cdot10^{-7}.
\]
Hence $4ky^4/(1-y)<4\,(25600/0.53)^2\,10^{-11}/(1-y)<0.09337$, and the second term in
Proposition~\ref{prop:majorant} equals $4\cdot80^4(y/40)^3/(1-80/k)<3.1\cdot10^{-7}$. Therefore
$\|h\|_\cA<0.0934$ and $\|e^h-1\|_\cA<e^{0.0934}-1<0.098$ (Appendix~\ref{app:cert}).

Let $H'=[b_{d-i-j}]_{i,j=0}^d$. Since $b_s=c_se^{h(s)}$ at the integers $|s|\le d$, $H'-M=M\circ(e^h-1)$, and
Corollary~\ref{cor:calculus} gives
$\|S^{-1}(H'-M)S^{-T}\|\le\|e^h-1\|_\cA<1$. By \eqref{eq:SJS} and Lemma~\ref{lem:inertia}, $H'$ has the inertia of $J$;
$J$ has $\lfloor r/2\rfloor$ negative entries, so $\sgn\det H'=(-1)^{\lfloor r/2\rfloor}$.

Finally, $a_{k+d-i-j}=a_k\rho^d\,\rho^{-i}\rho^{-j}\,b_{d-i-j}$ (all indices are $\ge k-d>0$), so
$H=[a_{k+d-i-j}]$ is a positive diagonal congruence of a positive multiple of $H'$, and $\sgn\det H=(-1)^{\lfloor r/2\rfloor}$.
Reversing the order of the columns of $[a_{k+j-i}]$ gives $H$, at the cost of the sign of the reversal permutation,
$(-1)^{\lfloor r/2\rfloor}$. Hence $D_{r,k}=(-1)^{\lfloor r/2\rfloor}\det H>0$.
\end{proof}

\begin{remark}[The constant]\label{rem:constant}
The main term is at most $4(25600/0.53)^2/C$ when $k\ge Cr^2$, so the argument closes as soon as this is $<\log2$, i.e.
for every $C>1.35\cdot10^{10}$. Using the cruder windows $k\theta>0.4$, $\alpha>4/3$ one gets $C>2.37\cdot10^{10}$. No effort
has been made to optimize Input~A, which is the dominant loss (the true coefficients satisfy
$|f^{(n)}(k)/n!|\,k^{n-1}<0.3$, cf.\ Table~\ref{tab:taylor}, versus $3\cdot40^n$).
\end{remark}

\section{Numerical illustration}\label{sec:numerics}

Nothing in this section is used in the proof. We use the certified values of $\lambda_n$, $n\le12016$, published
with \cite{LM} (relative width $<3\cdot10^{-91}$), which we cross-checked against our own quadrature of \eqref{eq:Xi}.
At $k=12000$ one finds $k\tau_k=1.74995$, $k^2|\tau_k-\cT(k)|=1.017$ (compare $20$ in Input~B), $k^2\Delta_k=1.7225$,
$\alpha=1.7777$, $k\theta=0.7655$.

\begin{table}[h]
\centering
\begin{tabular}{rrcccc}
\toprule
$k$ & $r$ & three-point & four-point & ratio & four-point$\cdot k/r^2$\\
\midrule
11990 & 3 & $4.86\cdot10^{-3}$ & $1.77\cdot10^{-5}$ & 275 & 0.024\\
11990 & 6 & $4.00\cdot10^{-2}$ & $2.17\cdot10^{-4}$ & 184 & 0.072\\
11990 & 8 & $7.93\cdot10^{-2}$ & $5.02\cdot10^{-4}$ & 158 & 0.094\\
11990 & 12 & $1.88\cdot10^{-1}$ & $1.49\cdot10^{-3}$ & 127 & 0.124\\
11990 & 16 & $3.31\cdot10^{-1}$ & $3.06\cdot10^{-3}$ & 108 & 0.144\\
\midrule
3000 & 8 & $1.61\cdot10^{-1}$ & $1.93\cdot10^{-3}$ & 83 & 0.091\\
6000 & 8 & $1.13\cdot10^{-1}$ & $9.86\cdot10^{-4}$ & 114 & 0.092\\
\bottomrule
\end{tabular}
\caption{Actual whitened perturbation $\|S^{-1}(H'-M)S^{-T}\|$ on the true $\xi$-coefficients, for the three-point
reference of \cite{Mich} and the four-point reference of this paper. At fixed $r$ the four-point error decays like $k^{-1}$
and the three-point error like $k^{-1/2}$.}
\label{tab:whitened}
\end{table}

\begin{table}[h]
\centering
\begin{tabular}{rcccc}
\toprule
$n$ & $|f^{(n)}(k)/n!|\,k^{n-1}$ & $|[s^n]\ell|\,k^{n-1}$ & $|\beta_n|\,k^{n-1}$ & bound $4\cdot40^n$\\
\midrule
3 & 0.287 & 0.287 & $9.0\cdot10^{-6}$ & $2.6\cdot10^{5}$\\
4 & 0.142 & 0.196 & 0.054 & $1.0\cdot10^{7}$\\
5 & 0.085 & 0.157 & 0.072 & $4.1\cdot10^{8}$\\
6 & 0.056 & 0.139 & 0.083 & $1.6\cdot10^{10}$\\
8 & 0.030 & 0.132 & 0.102 & $2.6\cdot10^{13}$\\
\bottomrule
\end{tabular}
\caption{Taylor coefficients at $k=12000$, computed by Cauchy integrals of $\log I(z)-\log\Gamma(2z+1)$ and of \eqref{eq:logc}.
The predicted size of the dominant term $\beta_4\cR(P_4)$ at $(k,r)=(11990,8)$ is $5.0\cdot10^{-4}$, in agreement with
Table~\ref{tab:whitened}.}
\label{tab:taylor}
\end{table}

\begin{remark}[Outlook]
Heuristically, a reference with $m$ exact matching points and an exact $LDL^T$ with $q$-Pascal factor would give a main
term $\asymp k^{1-m}w^m\asymp r^{m/2}k^{1-m/2}$, i.e.\ a wedge $k\gtrsim r^{m/(m-2)}$. Whether such references exist for
$m\ge5$ (for instance among ${}_r\phi_s$-type sequences) is open; in any case this route cannot reach the critical regime
$k=O(r)$, where $w$ becomes comparable to the radius of analyticity.
\end{remark}

\appendix
\section{Certified scalar inequalities and reproducibility}\label{app:cert}

All scalar inequalities in Sections~\ref{sec:windows}--\ref{sec:proof} are verified by the ancillary script
\texttt{certificate.py} using Arb ball arithmetic \cite{Arb} at 256 bits. Each window has the form $c(k)/k$ or
$c(k)/k^2$ with $c$ monotone in the favourable direction, so evaluation at $k=12000$ suffices. The certified values are:

\begin{center}
\begin{tabular}{lll}
\toprule
quantity & certified enclosure & used\\
\midrule
$k\tau_k$ & $[1.712619,\ 2.001667]$ & $(1.71,\,2.002)$\\
$k^2\Delta_k$ & $[1.675686,\ 2.003334]$ & $(1.675,\,2.004)$\\
$\alpha$ & $[1.458989,\ 2.392839]$ & $(1.45,\,2.40)$\\
$k\theta$ & $[0.537774,\ 2.002]$ & $(0.53,\,2.002)$\\
$k\kappa$ & $\le 6.451$ & $<8.4$\\
$y$ (wedge) & $\le4.915\cdot10^{-4}$ & \\
$\|h\|_\cA$ (wedge) & $\le0.0933691$ & $<0.0934$\\
$\|e^h-1\|_\cA$ & $<0.098$ & $<1$\\
\bottomrule
\end{tabular}
\end{center}

\noindent Further ancillary scripts (not needed for the proof) check:
\texttt{sym\_check.py} --- symbolic verification (generic parameters) of Lemma~\ref{lem:match}, the integer cocycle,
Lemma~\ref{lem:congruence}, Proposition~\ref{prop:LDL} for $i,j<7$, Lemmas~\ref{lem:R1}--\ref{lem:subdiag} and
Lemma~\ref{lem:interp};
\texttt{gateB\_numeric.py}, \texttt{response\_scaling.py} --- the response bound of Proposition~\ref{prop:response},
computed from a numerical $LDL^T$ of $M$ at 260 digits;
\texttt{xi\_data.py}, \texttt{whitened\_experiment.py}, \texttt{taylor\_check.py} --- Section~\ref{sec:numerics}.

\subsection*{Acknowledgments}
The author thanks W.~Michałowski and Y.~Liu and J.~Mao for making their preprints and certified data publicly available.

\subsection*{Use of AI tools}
The following AI models were used in the preparation of this work: Astra (GPT, OpenAI) and Claude (Anthropic).

\begin{thebibliography}{99}

\bibitem{ASW}
M.~Aissen, I.~J.~Schoenberg, A.~M.~Whitney,
\emph{On the generating functions of totally positive sequences I},
J.~Analyse Math.\ \textbf{2} (1952), 93--103.

\bibitem{CNV}
G.~Csordas, T.~S.~Norfolk, R.~S.~Varga,
\emph{The Riemann hypothesis and the Tur\'an inequalities},
Trans.\ Amer.\ Math.\ Soc.\ \textbf{296} (1986), 521--541, doi:10.2307/2000378.

\bibitem{Edrei}
A.~Edrei,
\emph{On the generating functions of totally positive sequences II},
J.~Analyse Math.\ \textbf{2} (1952), doi:10.1007/BF02786971.

\bibitem{GR}
G.~Gasper, M.~Rahman,
\emph{Basic Hypergeometric Series}, 2nd ed.,
Cambridge University Press, 2004.

\bibitem{GORZ}
M.~Griffin, K.~Ono, L.~Rolen, D.~Zagier,
\emph{Jensen polynomials for the Riemann zeta function and other sequences},
Proc.\ Natl.\ Acad.\ Sci.\ USA \textbf{116} (2019), no.~23, 11103--11110.

\bibitem{Holland}
J.~Holland,
\emph{A new hyperbolicity wedge and a joint semicircle limit for Jensen polynomials of Riemann's $\xi$-function},
arXiv:2608.08682 (2026).

\bibitem{Arb}
F.~Johansson,
\emph{Arb: efficient arbitrary-precision midpoint-radius interval arithmetic},
IEEE Trans.\ Comput.\ \textbf{66} (2017), no.~8, 1281--1292.

\bibitem{Katkova}
O.~M.~Katkova,
\emph{Multiple positivity and the Riemann zeta-function},
Comput.\ Methods Funct.\ Theory \textbf{7} (2007), no.~1, 13--31; arXiv:math/0505174.

\bibitem{LM}
Y.~Liu, J.~Mao,
\emph{Infinite log-concavity of the Taylor coefficients of the Riemann xi-function},
arXiv:2610.04972 (2026).

\bibitem{Mich}
W.~Michałowski,
\emph{An explicit uniform cubic wedge for consecutive Toeplitz minors of the Riemann xi coefficients},
arXiv:2607.16795 (2026).

\bibitem{Nuttall}
J.~Nuttall,
\emph{Wronskians, cumulants, and the Riemann hypothesis},
Constr.\ Approx.\ \textbf{38} (2013), 193--212, doi:10.1007/s00365-013-9200-8.

\bibitem{Sra}
S.~Sra,
\emph{On some positivity properties of the Riemann $\Xi$ function},
arXiv:2611.02944 (2026).

\bibitem{WY}
L.~X.~W.~Wang, N.~N.~Y.~Yang,
\emph{Laguerre inequalities and complete monotonicity for the Riemann Xi-function and the partition function},
Trans.\ Amer.\ Math.\ Soc.\ (2024).

\end{thebibliography}

\end{document}
